\subsection{Высказывание, предикат, импликация}

Наша речь состоит из фраз и предложений, которые чаще всего являются некоторыми \foc{утверждениями}.
Однако не любые утверждения интересны математикам~--- ценность представляют утверждения, которым можно приписать конкретное \foc{истинностное значение}.

\defi{defi:высказывание}
{Высказывание}{
    Утверждение, о котором можно однозначно сказать, что оно \foc{истинно} (<<True>>, <<T>>) или \foc{ложно} (<<False>>, <<F>>).
}[
    \item Предложение <<Москва~--- столица России>>~--- истинное высказывание, а предложение <<Радиус Земли меньше радиуса Луны>>~--- ложное высказывание.
    \item Равенство $6 \mul 7 = 42$~--- истинное высказывание, а $2 + 2 = 5$~--- ложное высказывание.
    \item Равенство $x + 2 = 5$ не является высказыванием (это неопределённое высказывание или предикат).
    \item Предложение <<Привет, как дела?>> не является высказыванием.
][
    Предложение, о котором невозможно однозначно сказать, истинно оно или ложно, высказыванием не является.
    Например, к высказываниям не относятся предложение <<Данное утверждение ложно>> и следующая пара предложений, ссылающихся друг на друга: <<Следующее утверждение ложно. Предыдущее утверждение истинно>>.
]


\subsection{Математические выражения и определения}

В математическом тексте внешне похожие записи могут играть существенно разные роли.
Важно различать три типа конструкций:
\foc{алгебраические выражения},
\foc{логические выражения} и
\foc{определения}.

Алгебраическое выражение обозначает некоторый математический объект или значение.
Примерами служат $42$, $x$, $2 + 2$, $x^2 - 3x + 2$ и~т.п.
Если значения всех входящих в выражение переменных зафиксированы, то выражение получает значение.
При этом само алгебраическое выражение не является ни истинным, ни ложным.
Например, при $x = 3$ выражение $x^2 - 3x + 2$ равно $2$, но не является высказыванием.

Логическое выражение, напротив, утверждает нечто об объектах и потому либо имеет истинностное значение само по себе, либо приобретает его после подстановки конкретных значений переменных.
Например, $2 + 2 = 4$ и $5 > 3$~--- высказывания, поскольку их истинность может быть установлена непосредственно.
Выражения $x + 2 = 5$ и $x^2 > 0$~--- предикаты: после подстановки конкретного значения $x$ каждое из них становится истинным или ложным высказыванием.

Знак равенства, как и знаки $<$, $>$, $\le$, $\ge$, $\ne$, превращает два алгебраических выражения в логическое выражение.
Например, $x + 2$~--- алгебраическое выражение, а $x + 2 = 5$~--- логическое выражение.
С этой точки зрения уравнение $x + 2 = 5$ является предикатом от переменной $x$.
Решить это уравнение означает найти все значения $x$, при которых данный предикат истинен.

Для введения новых обозначений мы используем специальные символы определения.
Если новое обозначение вводится для математического объекта или значения, применяется равенство по определению:
\begin{equation*}
    s \is a + b,
    \quad
    a + b \isr s.
\end{equation*}
Обе записи означают, что символ $s$ далее обозначает выражение $a + b$.
Двоеточие ставится со стороны нового обозначения.
Заметим, что в оборотах <<обозначим>>, <<положим>>, <<введём обозначение>> и~т.п. определяющий смысл часто ясен из контекста.
Поэтому для локальных вспомогательных величин допустима и более привычная запись через обычный знак $=$.
Однако при первом введении важных обозначений мы будем по возможности использовать явные символы $\is$ и $\isr$.

Если новое обозначение вводится для высказывания, предиката или свойства, используется знак равносильности по определению:
\begin{equation*}
    A \eqdef (2 + 2 = 4),
    \quad
    (x > 0) \eqdefr P(x).
\end{equation*}
Первая запись вводит имя $A$ для высказывания $2 + 2 = 4$, а вторая~--- имя $P(x)$ для предиката $x > 0$.

Символы $\is$, $\isr$, $\eqdef$, $\eqdefr$ не сравнивают уже заданные объекты, а вводят соглашение о новом обозначении.
Запись $s \is 2 + 2$ сама по себе не является истинным или ложным высказыванием~--- она лишь вводит обозначение $s$.
Когда обозначение введено, то, например, записи $s = 4$ и $s = 5$ уже являются обычными высказываниями: первое истинно, второе ложно.
Аналогично, после определения $P(x) \eqdef (x^2 > 0)$ запись $P(x) \eq (x \ne 0)$ является предикатом.
При действительных $x$ этот предикат истинен при всех $x$, поскольку условия $x^2 > 0$ и $x \ne 0$ равносильны.

Похожее различие возникает и в языках программирования.
Например, в Python оператор присваивания \mintinline{python}{=} и объявление функции с помощью ключевого слова \mintinline{python}{def} вводят имена для объектов или правил вычисления.
Они не являются проверками и не возвращают истинностное значение.
Оператор сравнения \mintinline{python}{==}, напротив, проверяет равенство двух уже заданных значений и возвращает \mintinline{python}{True} или \mintinline{python}{False}.
Эти различия иллюстрируются в листинге~\ref{code:выражения_и_определения_в_python}.

\code
    {code:выражения_и_определения_в_python}
    {Выражения и определения в Python}
    {code/l-01/code1.py}
    [code/l-01/code1.txt]